\begin{thebibliography}{99}
\bibitem{gama} J. Gama et al., ``A survey on concept drift adaptation'', \emph{ACM Comput. Surv.}, vol. 46, no. 4, art. 44, 2014, doi: 10.1145/2523813.
\bibitem{lu} J. Lu et al., ``Learning under Concept Drift: A Review'', \emph{IEEE TKDE}, vol. 31, no. 12, pp. 2346--2363, 2019, doi: 10.1109/TKDE.2018.2876857.
\bibitem{sharaf} I. Sharafaldin, A. H. Lashkari, and A. A. Ghorbani, ``Toward Generating a New Intrusion Detection Dataset and Intrusion Traffic Characterization'', \emph{ICISSP}, pp. 108--116, 2018, doi: 10.5220/0006639801080116.
\bibitem{engelen} G. Engelen, V. Rimmer, and W. Joosen, ``Troubleshooting an Intrusion Detection Dataset: the CICIDS2017 Case Study'', \emph{IEEE SPW/WTMC}, pp. 7--12, 2021, doi: 10.1109/SPW53761.2021.00009.
\bibitem{liu} L. Liu et al., ``Error Prevalence in NIDS datasets: A Case Study on CIC-IDS-2017 and CSE-CIC-IDS-2018'', \emph{IEEE CNS}, 2022.
\bibitem{arp} D. Arp et al., ``Dos and Don'ts of Machine Learning in Computer Security'', \emph{USENIX Security}, pp. 3971--3988, 2022.
\bibitem{tesseract} F. Pendlebury et al., ``TESSERACT: Eliminating Experimental Bias in Malware Classification across Space and Time'', \emph{USENIX Security}, pp. 729--746, 2019.
\bibitem{paired} S. H. Bach and M. A. Maloof, ``Paired Learners for Concept Drift'', \emph{IEEE ICDM}, pp. 23--32, 2008.
\bibitem{sam} V. Losing, B. Hammer, and H. Wersing, ``Self-Adjusting Memory: How to Deal with Diverse Drift Types'', \emph{IJCAI}, pp. 4899--4903, 2017, doi: 10.24963/ijcai.2017/690.
\bibitem{adwin} A. Bifet and R. Gavalda, ``Learning from Time-Changing Data with Adaptive Windowing'', \emph{SIAM SDM}, pp. 443--448, 2007, doi: 10.1137/1.9781611972771.42.
\bibitem{arf} H. M. Gomes et al., ``Adaptive random forests for evolving data stream classification'', \emph{Mach. Learn.}, vol. 106, pp. 1469--1495, 2017, doi: 10.1007/s10994-017-5642-8.
\bibitem{soul} S. K. Amalapuram et al., ``SOUL: A Semi-supervised Open-world continUal Learning method for Network Intrusion Detection'', \emph{arXiv:2412.00911}, 2024.
\bibitem{ssf} X. Zhang et al., ``Continual Learning with Strategic Selection and Forgetting for Network Intrusion Detection'', \emph{IEEE INFOCOM}, arXiv:2412.16264v4, 2025.
\bibitem{jeon} S. Jeon and A. Kim, ``A Label-Efficient Adaptive Retraining Framework for Network Intrusion Detection Under Streaming Arrival Bias'', \emph{Electronics}, vol. 15, no. 17, art. 3812, 2026, doi: 10.3390/electronics15173812.
\bibitem{nafn} Z. Zhao et al., ``Network Intrusion Detection Adaptability under Spurious Labels: A Noise-Aware Few-Shot Class-Incremental Learning Framework'', \emph{IEEE TNSM}, 2026, doi: 10.1109/TNSM.2026.3752751.
\bibitem{acorn} S. Fuhrman, O. Gungor, and T. Rosing, ``ACORN-IDS: Adaptive Continual Novelty Detection for Intrusion Detection Systems'', \emph{arXiv:2602.07291}, 2026.
\bibitem{friedman} J. H. Friedman, ``Greedy function approximation: A gradient boosting machine'', \emph{Ann. Statist.}, vol. 29, no. 5, pp. 1189--1232, 2001, doi: 10.1214/aos/1013203451.
\bibitem{lightgbm} G. Ke et al., ``LightGBM: A Highly Efficient Gradient Boosting Decision Tree'', \emph{NeurIPS}, vol. 30, pp. 3146--3154, 2017.
\bibitem{sklearn} scikit-learn developers, ``Histogram-Based Gradient Boosting'', \emph{scikit-learn 1.9.1 documentation}, accessed Nov. 26, 2026, \url{https://scikit-learn.org/stable/modules/ensemble.html}.
\bibitem{cover} T. Cover and P. Hart, ``Nearest neighbor pattern classification'', \emph{IEEE Trans. Inf. Theory}, vol. 13, no. 1, pp. 21--27, 1967, doi: 10.1109/TIT.1967.1053964.
\bibitem{dudani} S. A. Dudani, ``The Distance-Weighted k-Nearest-Neighbor Rule'', \emph{IEEE Trans. Syst., Man, Cybern.}, vol. SMC-6, no. 4, pp. 325--327, 1976, doi: 10.1109/TSMC.1976.5408784.
\bibitem{guo} C. Guo et al., ``On Calibration of Modern Neural Networks'', \emph{ICML/PMLR}, vol. 70, pp. 1321--1330, 2017.
\bibitem{settles} B. Settles, ``Active Learning Literature Survey'', \emph{Univ. Wisconsin--Madison}, Tech. Rep. 1648, revised Jan. 2010, 2009.
\bibitem{tavallaee} M. Tavallaee et al., ``A Detailed Analysis of the KDD CUP 99 Data Set'', \emph{IEEE CISDA}, pp. 1--6, 2009, doi: 10.1109/CISDA.2009.5356528.
\bibitem{moustafa} N. Moustafa and J. Slay, ``UNSW-NB15: a comprehensive data set for network intrusion detection systems (UNSW-NB15 network data set)'', \emph{IEEE MilCIS}, pp. 1--6, 2015, doi: 10.1109/MilCIS.2015.7348942.
\bibitem{river} River developers, ``ARFClassifier'', \emph{River 0.22.0 documentation}, accessed Nov. 26, 2026, \url{https://riverml.xyz/0.22.0/api/forest/ARFClassifier/}.
\bibitem{saito} T. Saito and M. Rehmsmeier, ``The Precision-Recall Plot Is More Informative than the ROC Plot When Evaluating Binary Classifiers on Imbalanced Datasets'', \emph{PLoS ONE}, vol. 10, no. 3, e0118432, 2015, doi: 10.1371/journal.pone.0118432.
\end{thebibliography}